\section{Introduction and theorem structure}\label{sec:current-overview}\label{sec:introduction}
Full-history boundary attainment, finite-memory approximation, and exact finite-memory attainment are different assertions. The performed-law theory gives the first under a positive call-cost floor, but not the third. The weighted boundary theory already derives an optimal probability-valued decision face and characterizes its exact stochastic implementation by an anchored invariant convex family. Here the central question is how far that classification and its physical consequences extend beyond its original sufficient stopping certificate.

Theorem~\ref{thm:r67-rigidity} supplies an exact stopping criterion. For each possible first-herald word, compare its probability with the unused budget in each model. Every boundary policy must stop there precisely when the infimum of this normalized slack is zero. The converse is constructive: positive slack at one word permits a finite-memory policy to reject there and remain within the same budget. This covers state-dependent heralds and model-dependent failures without replacing the actual report law by a reference execution. Example~\ref{ex:r67-flat-margin} shows that cost-active support coverage is genuinely stronger than necessary.

Theorems~\ref{thm:weighted-face} and~\ref{thm:convex-memory}, preserved in the supplement, already solve robust ties by positive realization. The pointed-polytope statement in Theorem~\ref{thm:r67-polytope} extends that existing conclusion to the exact rigidity class and records the cost of a prescribed initial label. Classical residual automata, linear representations and positive realization provide its mathematical language \cite{RabinScott59,Jaeger00,ThonJaeger15,MonrasWinter14}; a finite linear rank is not a positivity certificate. We do not claim to introduce these methods or to solve an arbitrary finite-memory POMDP by renaming a residual function.

The new physical realization concerns the same ordinary squared loss in every model. Theorem~\ref{thm:r67-embedding} uses a protected channel, one additional classical register and a three-outcome prospective audit to force replication of the channel's response probabilities. The optimal randomized kernel is unique even though an attaining least-favorable measure ties all actions at all words. This differs from the signed calibration loss of Theorem~\ref{thm:calibration-realization}. The Erlang-chain memory hierarchy of Theorem~\ref{thm:calibration-k} and the unit-circle obstruction of Theorem~\ref{thm:qubit-positive-phase} are inherited mechanisms; they are realized here by a common squared audit, not advertised as discoveries of this continuation.

Proposition~\ref{prop:r67-finite-witness} adds a further conclusion that does not require an interval of unknown models. A finite rational set of preparations and instruments already separates every smaller stochastic observer from the target memory order by an explicit positive risk margin. A finite response determinant gives the memory obstruction; inverse Vandermonde bounds and a controlled series tail turn it into finitely many actual model witnesses. No finite sampling argument is used in place of that proof. The resulting lower is deliberately not described as a sharp large-memory rate for the rotating qubit.

Theorem~\ref{thm:r67-algebraic-memory} optimizes the unknown robust risk itself at each fixed memory. In a compact semialgebraic instrument family, bounded score potentials on the joint physical--observer space yield a quantified polynomial system. Its attained endpoint is algebraic, and real-algebraic elimination supplies a controller. This differs from testing realization of an already specified rational response in Proposition~\ref{prop:finite-realization-test}. It does not decide whether some unspecified finite memory attains the unrestricted optimum, and it provides no efficient universal bit-complexity bound.

\subsection*{Logical and resource boundaries}
The new core assumes one globally legal command, a fixed finite dictionary, a compact continuous model family, a primitive uniformly transient first-herald instrument, bounded prospective loss and positive bounded call costs. The exact stopping criterion is an additional conclusion about the minimum-cost face, not a premise concealed in an audit. Controllers see only the declared reports and retain counted classical labels; the protected physical state and a least-favorable model measure are not runtime inputs. Event stopping packs live and decoder roles by a maximum, while state stopping keeps them disjoint. Program encodings and working storage are further resources.

The main article gives complete proofs of the new stopping, realization and finite-witness results, and retains the original actual-history geometric transfer proof as a complementary general theorem. Only the elementary determinant lemma in Supplement~\ref{sec:r67-determinant} is an additional technical input to the finite witness. The remaining supplement consists of independent current developments with their full proofs: performed-law compactness and duality, primitive and compact-model synthesis, the weighted realization and quantum phase theorems, timing, width and rank geometry, reset control, realizations, composition and fixed-law predictive representation. Purely historical source preservation is outside this minimal dependency set. No conclusion for nonintervening sensing, ordered measurements or two-stage reset communication is inferred from the squared-audit example without its own theorem.
